\documentclass[border=3pt,tikz]{standalone}
\usepackage{physics}
\usepackage{siunitx}
\usepackage{ifthen}
\usepackage{tikz}
\usepackage[outline]{contour} % glow around text
\usepackage{silence}

\usetikzlibrary{calc}
\AtBeginDocument{\RenewCommandCopy\qty\SI}
\usetikzlibrary{angles,quotes} % for pic
\usetikzlibrary{arrows.meta}
\usetikzlibrary{patterns}
\tikzset{>=latex} % for LaTeX arrow head
\contourlength{1.4pt}

% Reusable coordinate system pic
% Usage: \pic at (x,y) {coordsys};
%   or:  \pic[rotate=30] at (x,y) {coordsys};
\tikzset{
  coordsys/.pic={
    \fill[red] (0,0) circle(0.1);
    \draw[->,thick, red] (0,0) -- (1.4,0) node[below] {$x$};
    \draw[->,thick, red] (0,0) -- (0,1.4) node[above] {$y$};
  }
}

% Reusable potentiometer pic (concentric angular sensor)
% Expects \ang (pendulum angle from -y) and \potR (outer radius) to be defined.
\def\potR{0.85}       % potentiometer outer housing radius
\def\trackR{0.63}     % resistive track radius
\tikzset{
  potentiometer/.pic={
    % Outer housing
    \draw[thick, fill=white!92!black] (0,0) circle (\potR);
    % Resistive track (270° arc, dead zone at top: 45°–135°)
    \draw[line width=2.5pt, black!35]
        (135:\trackR) arc (135:-135:\trackR);
    % Track end caps
    \fill[black!45] ( 135:\trackR) circle (0.055);
    \fill[black!45] (-135:\trackR) circle (0.055);
    % Wiper arm (aligned to pendulum angle \ang from -y axis)
    \pgfmathsetmacro{\wiperAng}{\ang - 90}
    \draw[semithick, black!55] (0,0) -- (\wiperAng:\trackR);
    \fill[black!65] (\wiperAng:\trackR) circle (0.05);
  }
}

\colorlet{xcol}{blue!70!black}
\colorlet{vcol}{green!60!black}
\colorlet{myred}{red!65!black}
\colorlet{acol}{red!50!blue!80!black!80}
\tikzstyle{ground}=[preaction={fill,top color=black!10,bottom color=black!5,shading angle=20},
                    fill,pattern=north east lines,draw=none,minimum width=0.3,minimum height=0.6]
\tikzstyle{mass}=[line width=0.6,red!30!black,fill=red!40!black!10,rounded corners=1,
                  top color=red!40!black!20,bottom color=red!40!black!10,shading angle=20]
\tikzstyle{faded mass}=[dashed,line width=0.1,red!30!black!40,fill=red!40!black!10,rounded corners=1,
                        top color=red!40!black!10,bottom color=red!40!black!10,shading angle=20]
\tikzstyle{rod}=[brown!70!black,very thick,line cap=round]
\def\rod#1{ \draw[black,line width=1.4] #1; \draw[rod,line width=1.1] #1; }
\tikzstyle{force}=[->,myred,very thick,line cap=round]
\tikzstyle{velocity}=[->,vcol,very thick,line cap=round]

\begin{document}


% PENDULUM - ORIGIN AT HINGE
\def\L{4}      % rod length
\def\W{0.3}    % rod width
\def\R{0.4}    % pendulum head radius
\def\ang{40}   % angle from negative y-axis (hanging down)
\def\F{1.6}    % force arrow magnitude

\begin{tikzpicture}
  \coordinate (O) at (0,0);
  
  % Pendulum rod (rectangle, rotated about hinge)
  \begin{scope}[rotate=\ang]
    \draw[thick, fill=brown!30] (-\W/2,0) rectangle (\W/2,-\L);
    % Pendulum head (circle at end of rod)
    \draw[thick, fill=gray!40] (0,-\L) circle(\R);
    \node[right=\R+12pt] at (0,-\L) {$m,\,I$};
    % Rod midline (dash-dot) — extended beyond head for force reference
    \draw[dash dot, thin] (0,0.3) -- (0,-\L-\R-1.2);
    % Pendulum head center lines (dash-dot, cross at center of circle)
    \draw[dash dot, thin] (-\R-0.3,-\L) -- (\R+0.3,-\L);
    % Origin center line parallel to pendulum head center line (dash-dot)
    \draw[dash dot, thin] (-\R-0.3,0) -- (\R+0.3,0);
  \end{scope}
  
  % === Ghost pendulum at theta = 0 (contact position) ===
  \begin{scope}[opacity=0.35]
    \draw[thick, fill=brown!15] (-\W/2,0) rectangle (\W/2,-\L);
    \draw[thick, fill=gray!20] (0,-\L) circle(\R);
  \end{scope}
  
  % Vertical reference line (theta = 0)
  \draw[dash dot, thin] (0,0) -- (0,-\L-\R-0.3);
  
  % Pendulum trajectory (arc of the head center)
  \draw[dotted, thick] (-90:\L) arc (-90:-90+\ang+0:\L);
  
  % Wall (impact surface at theta = 0, right edge at x = -R)
  \def\wallW{0.4}   % wall thickness
  \def\wallH{2}   % wall height
  \draw[ground] (-\R-\wallW, -\L-\wallH/2) rectangle (-\R, -\L+\wallH/2);
  \draw[thick] (-\R, -\L-\wallH/2) -- (-\R, -\L+\wallH/2);
  % Wall label (with leader line)
  \node[left, font=\footnotesize\itshape] (walllabel) at (-\R-\wallW-0.8, -\L) {rigid wall};
  \draw[thin, gray] (walllabel.east) -- (-\R-\wallW, -\L);
  
  % Pendulum head center coordinate (needed for angle annotation)
  \coordinate (M) at (\ang-90:\L);
  
  % Angle theta at hinge (between vertical and rod midline)
  \coordinate (B) at (0,-1.5);
  \draw pic[->,"\,$\theta$", draw, angle radius=60, angle eccentricity=1.1] {angle=B--O--M};
  
  % Length L annotation (offset so it clears the hinge circle)
  \draw[<->,thick] (\ang-90+90:0.7) ++ (\ang-90:0) -- ++ (\ang-90:\L)
    node[midway, above, sloped] {$L$};
  
  % === Potentiometer (angular sensor) concentric with hinge ===
  \pic at (0,0) {potentiometer};
  
  % Terminal wires from potentiometer housing (left side, horizontal)
  % +5V supply (red) — exits at 150° on housing
  \draw[semithick, red!70!black]
      (150:\potR) -- (-2.2,{\potR*sin(150)});
  \fill[red!70!black] (-2.2,{\potR*sin(150)}) circle(0.05);
  \node[left, font=\footnotesize, red!70!black]
      at (-2.2,{\potR*sin(150)}) {$+5\,\mathrm{V}$};
  % Signal output V(θ) (green) — exits at 180° on housing
  \draw[semithick, green!50!black]
      (180:\potR) -- (-2.2,0);
  \fill[green!50!black] (-2.2,0) circle(0.05);
  \node[left, font=\footnotesize, green!50!black]
      at (-2.2,0) {$V(\theta)$};
  % Ground reference (black) — exits at 210° on housing
  \draw[semithick, black!70]
      (210:\potR) -- (-2.2,{\potR*sin(210)});
  \fill[black!70] (-2.2,{\potR*sin(210)}) circle(0.05);
  \node[left, font=\footnotesize, black!70]
      at (-2.2,{\potR*sin(210)}) {GND};
  % Sensor label
  \node[font=\footnotesize\itshape, text=black!60] at (-1.5,-0.85) {angular sensor};
  
  % === BLDC Motor with belt drive ===
  \def\motorX{3.5}       % motor shaft x-position
  \def\motorY{1.5}       % motor shaft y-position
  \def\motorW{1.4}       % motor body width (extends right from shaft)
  \def\motorH{1.0}       % motor body height
  \def\motorShaftR{0.35} % motor shaft/drive pulley radius
  
  % Belt tangent geometry (open belt: potR as driven pulley, motorShaftR as drive pulley)
  \pgfmathsetmacro{\beltAngle}{atan2(\motorY, \motorX)}
  \pgfmathsetmacro{\beltDist}{sqrt(\motorX*\motorX + \motorY*\motorY)}
  \pgfmathsetmacro{\beltBeta}{asin((\potR - \motorShaftR) / \beltDist)}
  \pgfmathsetmacro{\tangTop}{\beltAngle + 90 + \beltBeta}
  \pgfmathsetmacro{\tangBot}{\beltAngle - 90 - \beltBeta}
  % Belt lines (external tangent)
  \draw[semithick, black!45]
      (\tangTop:\potR) -- ($(\motorX,\motorY) + (\tangTop:\motorShaftR)$);
  \draw[semithick, black!45]
      (\tangBot:\potR) -- ($(\motorX,\motorY) + (\tangBot:\motorShaftR)$);
  
  % Motor body (extends right from shaft center)
  \draw[thick, fill=blue!5, rounded corners=2]
      (\motorX, \motorY-\motorH/2) rectangle (\motorX+\motorW, \motorY+\motorH/2);
  \node[font=\footnotesize] at (\motorX+\motorW/2+0.15, \motorY) {BLDC};
  % Motor shaft / drive pulley
  \draw[thick, fill=white] (\motorX, \motorY) circle(\motorShaftR);
  \fill[black!50] (\motorX, \motorY) circle(0.06);
  
  % Hinge joint / shaft center (drawn on top of potentiometer)
  \draw[thick, fill=white] (O) circle(0.4);
  % Revolute joint label (upper-left, with leader line)
  \node[font=\footnotesize\itshape] (jointlabel) at (-2.5,1.2) {revolute joint};
  \draw[thin, gray] (jointlabel.south east) -- (O);
  
  % Drive torque at hinge
  \draw[<-, very thick, blue!70!black] (160:1.0) arc (160:70:1.0)
    node[midway, above] {$\tau$};
  
  % Coordinate system at hinge (topmost layer)
  \pic at (0,0) {coordsys};
\end{tikzpicture}

\end{document}

